\writeSectionFrame{\presentationTitle}{Mediator}
\begin{frame}{Steckbrief}
	\begin{itemize}
 		\item Deutscher Name: Vermittler
 		\item Alternativer Name: Mediator
 		\item Englischer Name: Mediator
 		\item Gruppe: Verhaltensmuster
 	\end{itemize}
\end{frame}

\begin{frame}{Mediator}
	\begin{block}{Zweck}
 		Definiere ein Objekt, welches das Zusammenspiel einer Menge von Objekten in sich kapselt. Vermittler fördern lose Kopplung, indem sie Objekte davon abhalten, aufeinander explizit Bezug zu nehmen. Sie ermöglichen es, das Zusammenspiel der Objekte von ihnen unabhängig zu variieren.
 	\end{block}
\end{frame}

\begin{frame}{Allgemeines Klassendiagramm}
    \begin{tikzpicture}[scale=0.7, transform shape]
            \umlclass[type=interface, x=0, y=0]{Mediator}{}{
            }

            \umlclass[x=0, y=-4]{ConcreteMediator}{}
            {
                + addColleague(colleague: Colleague)\\
            }

            \umlclass[type=interface, x=8, y=0]{Colleague}{}{
                + setMediator(mediator: Mediator)\\
            }

            \umlclass[x=6, y=-4]{ConcreteColleague1}{}{
            }{
            }
        
            \umlclass[x=10, y=-4]{ConcreteColleague2}{}{
            }{
            }

            \umlimpl{ConcreteMediator}{Mediator}
            \umlimpl{ConcreteColleague1}{Colleague}
            \umlimpl{ConcreteColleague2}{Colleague}
            \umluniassoc{Colleague}{Mediator}
            \umluniassoc{ConcreteMediator}{Colleague}
    \end{tikzpicture}
\end{frame}

\begin{frame}{Telefonanlage}
	\begin{center}
 		\includegraphics[scale=0.3]{\BILDERPFAD/mediator1.jpg}
 	\end{center}
\end{frame}

\note{
	Wir möchten im ersten Beispiel gerne eine VOIP-Telefonanlage simulieren. Es gibt 
	verschiedene Arten von Teilnehmern. Jederzeit können neue Teilnehmer
	an den Switch angeschlossen und bereits existierende entfernt werden. VOIP-Geräte
	sind weitestgehend autonom, können z.B. selbstständig Verbindungen aufbauen. Der Switch
	erledigt Verwaltungsarbeit wie z.B. das Ermitteln eines Teilnehmers anhand der Telefonnummer
	oder die Erstellung von Protokollen oder Konferenzen. 
}

\begin{frame}{Klassendiagramm -- Telefonanlage}
    \begin{tikzpicture}[scale=0.7, transform shape]
            \umlclass[type=interface, x=0, y=0]{Switch}{}{
            }

            \umlclass[x=0, y=-4]{Telefonanlage}{}
            {
            }

            \umlclass[type=interface, x=8, y=0]{Teilnehmer}{}{
            }

            \umlclass[x=7, y=-5]{Fax}{}{
            }{
            }
        
            \umlclass[x=10, y=-6]{Telefon}{}{
            }{
            }

            \umlclass[x=4, y=-4]{PC}{}{
            }{
            }

            \umlimpl{Telefonanlage}{Switch}
            \umlimpl{Fax}{Teilnehmer}
            \umlimpl{PC}{Teilnehmer}
            \umlimpl{Telefon}{Teilnehmer}
            \umluniassoc{Teilnehmer}{Telefonanlage}
            \umluniassoc{Telefonanlage}{Fax}
            \umluniassoc{Telefonanlage}{PC}
            \umluniassoc{Telefonanlage}{Telefon}
    \end{tikzpicture}
\end{frame}

\begin{frame}[fragile]{Vermittler}
	\begin{exampleblock}{Switch}
 		\begin{lstlisting}
public interface Switch {
	void verbinden(Teilnehmer teilnehmer, 
			Teilnehmer gegenstelle);

	void verbindungTrennen(Teilnehmer teilnehmer, 
			Teilnehmer gegenstelle);

	void konferenzEroeffnen(Teilnehmer teilnehmer, 
			Teilnehmer[] gegenstellen);
}

        \end{lstlisting}
 	\end{exampleblock}
\end{frame}

\begin{frame}[fragile]{Konkreter Vermittler}
	\begin{exampleblock}{Telefonanlage}
 		\begin{lstlisting}
public class TelefonAnlage implements Switch {
	@Override
	public void verbinden(Teilnehmer teilnehmer, 
			Teilnehmer gegenstelle) {
		if (gegenstelle.istBesetzt()) {
			throw new SwitchException("...");
		}
		teilnehmer.verbinden(gegenstelle);
	}
	
	@Override
	public void konferenzEroeffnen(Teilnehmer teilnehmer, 
			Teilnehmer [] gegenstellen) {
		for (Teilnehmer gegenstelle: gegenstellen) {
			verbinden(teilnehmer, gegenstelle);
		}
	}
        \end{lstlisting}
 	\end{exampleblock}
\end{frame}

\begin{frame}[fragile]{Konkreter Vermittler}
	\begin{exampleblock}{Telefonanlage}
 		\begin{lstlisting}
	@Override
	public void verbindungTrennen(Teilnehmer teilnehmer, 
			Teilnehmer gegenstelle) {
		if (!teilnehmer.istBesetzt()) {
			throw new SwitchException("...");
		} else {
			teilnehmer.auflegen();
		}
		
		if (gegenstelle.istBesetzt()) {
			gegenstelle.auflegen();
		} 
	}
        \end{lstlisting}
 	\end{exampleblock}
\end{frame}

\begin{frame}[fragile]{Kollege}
	\begin{exampleblock}{Teilnehmer}
 		\begin{lstlisting}
public abstract class Teilnehmer {
	private Switch vermittler;
	private int rufnummer;
	private Teilnehmer aktuelleGegenstelle;
	
	public Teilnehmer(Switch vermittler, int rufnummer) {
		super();
        // ...
		aktuelleGegenstelle = null;
	}
	
	public boolean istBesetzt() {
		if (aktuelleGegenstelle == null) {
			return false;
		} else {
			return true;
		}
	}
        \end{lstlisting}
 	\end{exampleblock}
\end{frame}

\begin{frame}[fragile]{Kollege}
	\begin{exampleblock}{Teilnehmer}
 		\begin{lstlisting}
	public void anwaehlen(Teilnehmer gegenstelle) {
		try {
			vermittler.verbinden(this, gegenstelle);
		} catch (SwitchException s) {
			System.out.println(s.getMessage());
		}
	}
	
	public void verbinden(Teilnehmer gegenstelle) {
		aktuelleGegenstelle = gegenstelle;
		aktuelleGegenstelle.verbindungAnnehmen(this);
	}
	
	public void verbindungAnnehmen(Teilnehmer gegenstelle) {
		aktuelleGegenstelle = gegenstelle;
	}
        \end{lstlisting}
 	\end{exampleblock}
\end{frame}

\begin{frame}[fragile]{Kollege}
	\begin{exampleblock}{Teilnehmer}
 		\begin{lstlisting}
	public void auflegen() {
		aktuelleGegenstelle = null;
	}
	
	public void verbindungTrennen() {
		if (aktuelleGegenstelle != null) {
			vermittler.verbindungTrennen(this, aktuelleGegenstelle);
		}
	}
        \end{lstlisting}
 	\end{exampleblock}
\end{frame}

\begin{frame}[fragile]{Konkreter Kollege}
	\begin{exampleblock}{Telefon}
 		\begin{lstlisting}
public class Telefon extends Teilnehmer {

	public Telefon(Switch vermittler, int rufnummer) {
		super(vermittler, rufnummer);
	}
	
	public void anwaehlen(Teilnehmer[] gegenstellen) {
		try {
			getVermittler().konferenzEroeffnen(this, gegenstellen);
		} catch (SwitchException s) {
			System.out.println(s.getMessage());
		}
	}
}

        \end{lstlisting}
 	\end{exampleblock}
\end{frame}

\frameWithImageOnly{\BILDERPFAD/mediator2.jpg}
\note{
	Wir möchten einen kleinen Flughafen simulieren. In diesem Fall haben wir auch schon einen
	Mediator ins System integriert: den Fluglotsen, der auch in der Realität alle Anfragen
	bearbeitet und die Flugzeuge koordiniert. Dabei muss er darauf achten, dass er Flugzeuge
	Landebahnen zuweist, während Helikopter zu Heliports geschickt werden.
}

\begin{frame}{Klassendiagramm -- Flughafen}
    \begin{tikzpicture}[scale=0.7, transform shape]
        % Mediator Interface
        \umlinterface[x=-3,y=0]{Mediator}{}{}
    
        % FlightController Class
        \umlclass[x=-3, y=-3]{FlightController}{}{}
    
        % FlyingMachine Class
        \umlclass[x=3, y=0]{FlyingMachine}{}{}
    
        % Airplane Class
        \umlclass[x=1, y=-4]{Airplane}{}{}
    
        % Helicopter Class
        \umlclass[x=5, y=-4]{Helicopter}{}{}
    
        % LandingPlace Class
        \umlclass[x=-3, y=-6]{LandingPlace}{}{}
    
        % AirplaneRunway Class
        \umlclass[x=-5, y=-8]{AirplaneRunway}{}{}
    
        % Heliport Class
        \umlclass[x=-1, y=-8]{Heliport}{}{}
    
        % Relations
        \umlimpl{FlightController}{Mediator}
        \umluniassoc{FlightController}{FlyingMachine}
        \umluniassoc[anchor1=south, anchor2=north]{FlightController}{LandingPlace}
        \umlinherit[anchor1=north, anchor2=south]{Airplane}{FlyingMachine}
        \umlinherit[anchor1=north, anchor2=south]{Helicopter}{FlyingMachine}
        \umlinherit[anchor1=north, anchor2=south]{AirplaneRunway}{LandingPlace}
        \umlinherit[anchor1=north, anchor2=south]{Heliport}{LandingPlace}
    
    \end{tikzpicture}
\end{frame}

\begin{frame}[fragile]{Konkreter Vermittler}
	\begin{exampleblock}{FlightController}
 		\begin{lstlisting}
public class FlightController implements Mediator {
	private Map<String, FlyingMachine> flyingMachines;
	private List<LandingPlace> knownLandingPlaces;
	
	public void addLandingPlace(LandingPlace landingPlace) {
		knownLandingPlaces.add(landingPlace);
	}
	
	public void flyingMachineAppearsOnScreen(
            FlyingMachine flyingMachine) {
		flyingMachines.put(flyingMachine.
            getCallSignal(), flyingMachine);
	}
        \end{lstlisting}
 	\end{exampleblock}
\end{frame}

\begin{frame}[fragile]{Konkreter Vermittler}
	\begin{exampleblock}{FlightController}
 		\begin{lstlisting}
public void askLandingPermission(FlyingMachine flyingMachine) {
    if (flyingMachine instanceof Airplane) {
        LandingPlace landingPlace = findFreeAirplaneRunway();
        if (landingPlace != null) {
            landingPlace.land(flyingMachine);
        } else  {
            System.out.println("...");
        }
    } else if (flyingMachine instanceof Helicopter) {
        LandingPlace landingPlace = findFreeHeliport();
        if (landingPlace != null) {
            landingPlace.land(flyingMachine);
        } else  {
            System.out.println("...");
        }
    }
}
        \end{lstlisting}
 	\end{exampleblock}
\end{frame}

\begin{frame}[fragile]{Konkreter Vermittler}
	\begin{exampleblock}{FlightController}
 		\begin{lstlisting}
public void flyingMachineDisappearsFromScreen(
        FlyingMachine flyingMachine) {
    flyingMachines.remove(flyingMachine.
        getCallSignal(), flyingMachine);
}

public void flightMachineWantsToStart(FlyingMachine 
        flyingMachine) {
    for (LandingPlace landingPlace: knownLandingPlaces) {
        if (landingPlace.getLandedFlyingMachine().
                equals(flyingMachine)) {
            landingPlace.start(flyingMachine);
        }
    }
}
    \end{lstlisting}
 	\end{exampleblock}
\end{frame}

\begin{frame}[fragile]{Kollege}
	\begin{exampleblock}{FlyingMachine}
 		\begin{lstlisting}
public abstract class FlyingMachine {
	private String callSignal;
	private FlightController flightController;
	
	public FlyingMachine(String callSignal, 
			FlightController flightController) {
		this.callSignal = callSignal;
		this.flightController = flightController;
		flightController.flyingMachineAppearsOnScreen(this);
	}
	
	public String getCallSignal() {
		return callSignal;
	}
        \end{lstlisting}
 	\end{exampleblock}
\end{frame}

\begin{frame}[fragile]{Kollege}
	\begin{exampleblock}{FlyingMachine}
 		\begin{lstlisting}
public void askingForLandingPermission() {
    // Landevorbereitungen in der Maschine treffen.
    flightController.askLandingPermission(this);
}

public void iWantToStart() {
    // Treibstoff ueberpruefen usw.
    flightController.flightMachineWantsToStart(this);
}
        \end{lstlisting}
 	\end{exampleblock}
\end{frame}

\begin{frame}[fragile]{Konkreter Kollege}
	\begin{exampleblock}{Airplane}
 		\begin{lstlisting}
public class Airplane extends FlyingMachine {
	public Airplane(String callSignal, FlightController 
			flightController) {
		super(callSignal, flightController);
	}

	@Override
	public String toString() {
		return "Flugzeug: " + super.toString();
	}
}
        \end{lstlisting}
 	\end{exampleblock}
\end{frame}

\begin{frame}[fragile]{Konkreter Kollege 2}
	\begin{exampleblock}{LandingPlace}
 		\begin{lstlisting}
public class LandingPlace {
	private FlyingMachine landedFlyingMachine;
	private int number;
	
	public LandingPlace(int number) {
		this.number = number;
	}
	
	protected void land(FlyingMachine flyingMachine) {
		this.landedFlyingMachine = flyingMachine;
	}
	
	public void start(FlyingMachine flyingMachine) {
		this.landedFlyingMachine = null;
	}
        \end{lstlisting}
 	\end{exampleblock}
\end{frame}

\begin{frame}[fragile]{Konkreter Kollege 2}
	\begin{exampleblock}{AirplaneRunway}
 		\begin{lstlisting}
public boolean isFree() {
    if (landedFlyingMachine == null) {
        return true;
    } else {
        return false;
    }
}

public FlyingMachine getLandedFlyingMachine() {
    return landedFlyingMachine;
}
	
        \end{lstlisting}
 	\end{exampleblock}
\end{frame}

\frameWithImageOnly{\BILDERPFAD/mediator3.jpg}
\note{
	Und als letztes Beispiel sehen wir uns das Beispiel aus der Einführung an: die Verkehrskreuzung.
}

\begin{frame}{Klassendiagramm -- Verkehrskreuzung}
    \begin{tikzpicture}[scale=0.7, transform shape]
        \umlinterface[x=0, y=3]{Mediator}{}{}

        \umlclass[x=0, y=0]{TrafficMediator}{}{}
    
        \umlinterface[x=6, y=2]{CrossroadsComponent}{}{}
    
        \umlclass[x=8, y=-1]{TrafficLight}{}{}
    
        \umlclass[x=0, y=-3]{RailwayCrossing}{}{}
    
        \umlclass[x=4, y=-2]{Car}{}{}
    
        \umlimpl{TrafficMediator}{Mediator}
        \umlimpl{TrafficLight}{CrossroadsComponent}
        \umlimpl{RailwayCrossing}{CrossroadsComponent}
        \umlimpl{Car}{CrossroadsComponent}
        \umluniassoc[anchor1=west, anchor2=east, name=assoc1]{TrafficLight}{TrafficMediator}
        \umluniassoc{RailwayCrossing}{TrafficMediator}
        \umluniassoc[anchor1=west, anchor2=east, name=assoc3]{Car}{TrafficMediator}
    
    \end{tikzpicture}
\end{frame}

\begin{frame}[fragile]{Konkreter Vermittler}
	\begin{exampleblock}{TrafficMediator}
 		\begin{lstlisting}
public class TrafficMediator implements Mediator {
	private TrafficLight carTrafficLight;
	private TrafficLight pedestrianTrafficLight;
	private RailwayCrossing railwayCrossing;
	private List<Car> cars = new ArrayList<>();

	public void setCarTrafficLight(TrafficLight carTrafficLight) {
		this.carTrafficLight = carTrafficLight;
	}

	public void setPedestrianTrafficLight(
            TrafficLight pedestrianTrafficLight) {
		this.pedestrianTrafficLight = pedestrianTrafficLight;
	}
        \end{lstlisting}
 	\end{exampleblock}
\end{frame}

\begin{frame}[fragile]{Konkreter Vermittler}
	\begin{exampleblock}{TrafficMediator}
 		\begin{lstlisting}
@Override
public void notify(CrossroadsComponent sender, 
        CrossingEvent event) {
    if (sender == railwayCrossing &&
            event.equals(CrossingEvent.TRAIN_APPROACHING)) {
        carTrafficLight.setRed();
        pedestrianTrafficLight.setRed();
        for (Car car : cars) {
            car.stop();
        }
    } else if (sender == railwayCrossing && 
            event.equals(CrossingEvent.TRAIN_LEFT)) {
        carTrafficLight.setGreen();
        pedestrianTrafficLight.setGreen();
    } 
    // ...
}
        \end{lstlisting}
 	\end{exampleblock}
\end{frame}

\begin{frame}[fragile]{Kollege}
	\begin{exampleblock}{CrossroadsComponent}
 		\begin{lstlisting}
public interface CrossroadsComponent {
	void setMediator(Mediator mediator);
}
        \end{lstlisting}
 	\end{exampleblock}
\end{frame}

\begin{frame}[fragile]{Konkreter Kollege 1}
	\begin{exampleblock}{Car}
 		\begin{lstlisting}
public class Car implements CrossroadsComponent {
	private Mediator mediator;
	private CrossingEvent state;

	public Car() {
		this.state = CrossingEvent.CAR_STOPPED;
	}

	public void move() {
		state = CrossingEvent.CAR_MOVING;
		mediator.notify(this, state);
	}

	public void stop() {
		state = CrossingEvent.CAR_STOPPED;
		mediator.notify(this, state);
	}
}
        \end{lstlisting}
 	\end{exampleblock}
\end{frame}

\begin{frame}[fragile]{Konkreter Kollege 2}
	\begin{exampleblock}{TrafficLight}
 		\begin{lstlisting}
public class TrafficLight implements CrossroadsComponent {
	private Mediator mediator;
	private CrossingEvent state;

	public TrafficLight() {
		this.state = CrossingEvent.TRAFFIC_LIGHT_RED;
	}

	public void setRed() {
		state = CrossingEvent.TRAFFIC_LIGHT_RED;
		mediator.notify(this, state);
	}

	public void setGreen() {
		state = CrossingEvent.TRAFFIC_LIGHT_GREEN;
		mediator.notify(this, state);
	}
}
        \end{lstlisting}
 	\end{exampleblock}
\end{frame}
